\begin{proposition}\label{prop:decic111-excluded}
No elliptic curve over a number field of degree $10$ has a point
of order $111$.
\end{proposition}
\begin{proof}
Put $X=X_1(111)$ and $h=\langle10\rangle$.
We verify the hypotheses of Lemma~\ref{lem:dyadic-cuspidal-diamond}
with $d=10$.

There is no elliptic curve over $\mathbb F_{2^f}$ with order
divisible by $111$ for $1\leq f\leq10$.
For $f\leq6$ this follows from the Hasse bound.
The remaining Hasse-allowed possibilities are
\[
\begin{array}{c|c|c}
f&\#E(\mathbb F_{2^f})&a_{2^f}\\\hline
7&111&18\\
8&\text{none}&\text{none}\\
9&555&-42\\
10&999&26
\end{array}
\]
In characteristic $2$, an ordinary trace is odd.
The supersingular traces for odd $f$ are $0$ and
$\pm2^{(f+1)/2}$, and for even $f$ they are
$0$, $\pm2^{f/2}$ and $\pm2^{f/2+1}$;
see \cite[Theorem~4.1]{Waterhouse}.
None of the displayed even traces is allowed.
Thus $Y_1(111)(\mathbb F_{2^f})=\emptyset$ for $f\leq10$.

The global cusp signature is $1^{36},2^{18},36^2$.
To see this directly, write a Tate label as
$Q=\zeta_{111}^u q^{v/111}$, with $(u,v,111)=1$,
and quotient by $(u,v)\mapsto(u+kv,v)$ and simultaneous negation.
The cases $\gcd(v,111)=1,3,37,111$ give respectively
$36$ rational cusps, $18$ cusps over $\mathbb Q(\zeta_3)$,
one cusp over $\mathbb Q(\zeta_{37})$, and one cusp over
$\mathbb Q(\zeta_{111})^+$.
Their degrees are unchanged at $2$:
\[
 \operatorname{ord}_{3}(2)=2,\qquad
 \operatorname{ord}_{37}(2)=36,\qquad
 \operatorname{ord}_{(\mathbb Z/111\mathbb Z)^\times/\{\pm1\}}(2)=36.
\]
Hence every cuspidal divisor over $\mathbb F_2$ of degree $10$
lifts to a rational cuspidal divisor of degree $10$.

The projective diamond group is cyclic of order $36$, generated
by $\langle2\rangle$, and $10\equiv2^{24}\pmod {111}$.
A character is nontrivial on $h$ precisely when its order is
$9$, $18$ or $36$.
The complete exact winding computation gives
\[
\begin{array}{c|c|c|c}
\operatorname{ord}(\chi)&\operatorname{cond}(\chi)&
\dim S_\chi&\dim S_\chi^{\mathrm{wind}}\\\hline
9&37&11&11\\
18&37&11&11\\
36&111&11&11
\end{array}
\]
The full character list has weighted dimension $385=g(X)$.
For the order-$36$ component an exported $11\times11$ modular
Krylov matrix has rank $11$ at the good prime $73$.
The other two full winding dimensions are exact characteristic-zero
computations. It follows that $(h-1)J_1(111)(\mathbb Q)$ is finite.

Finally, $h$ has order $3$, and Abramovich's bound gives
\[
 \operatorname{gon}_{\mathbb C}(X)>
 \frac{325\cdot5472}{32768}
 =\frac{55575}{1024}>40.
\]
Since $3\nmid10$, the lemma applies.
\end{proof}
